\documentclass[11pt,a4paper]{article}

\usepackage[a4paper,margin=27mm,headsep=8mm,footskip=13mm]{geometry}
\IfFileExists{libertinus.sty}{\usepackage{libertinus}}{\usepackage{lmodern}}
\usepackage[T1]{fontenc}
\usepackage[utf8]{inputenc}
\usepackage{microtype}
\usepackage{amsmath,amssymb,amsthm,mathtools}
\usepackage{mathrsfs}
\usepackage{bm}
\usepackage{booktabs,longtable,tabularx,array}
\usepackage{enumitem}
\usepackage{xcolor}
\usepackage[most]{tcolorbox}
\usepackage{fancyhdr}
\usepackage{hyperref}
\usepackage[nameinlink,capitalise,noabbrev]{cleveref}
\usepackage{listings}
\usepackage{needspace}

\definecolor{deepolive}{HTML}{41533B}
\definecolor{softolive}{HTML}{F0F4EC}
\definecolor{deepteal}{HTML}{155E63}
\definecolor{softteal}{HTML}{EAF5F5}
\definecolor{warmgray}{HTML}{6A625C}
\definecolor{softcoral}{HTML}{FCF1EC}
\definecolor{deepblue}{HTML}{254D73}
\definecolor{softblue}{HTML}{EEF4FA}

\hypersetup{
  colorlinks=true,
  linkcolor=deepteal,
  citecolor=deepolive,
  urlcolor=deepblue,
  pdftitle={The Full Asymptotic Transseries of the Bell Numbers},
  pdfauthor={A repository-guided research synthesis}
}

\pagestyle{fancy}
\fancyhf{}
\fancyhead[L]{\small\textsc{Bell-number asymptotic transseries}}
\fancyhead[R]{\small\textsc{Coefficient calculus and Lambert--W geometry}}
\fancyfoot[C]{\thepage}
\renewcommand{\headrulewidth}{0.4pt}

\setlength{\parindent}{0pt}
\setlength{\parskip}{5.5pt plus 1.5pt minus 1pt}
\setcounter{tocdepth}{2}
\setcounter{secnumdepth}{3}
\allowdisplaybreaks
\emergencystretch=2.5em

\newtheorem{theorem}{Theorem}[section]
\newtheorem{lemma}[theorem]{Lemma}
\newtheorem{proposition}[theorem]{Proposition}
\newtheorem{corollary}[theorem]{Corollary}
\newtheorem{identity}[theorem]{Identity}
\newtheorem{conjecture}[theorem]{Conjecture}
\theoremstyle{definition}
\newtheorem{definition}[theorem]{Definition}
\newtheorem{algorithm}[theorem]{Algorithm}
\newtheorem{example}[theorem]{Example}
\theoremstyle{remark}
\newtheorem{remark}[theorem]{Remark}
\newtheorem{historical}[theorem]{Historical note}

\newtcolorbox{masterbox}[1][]{
  enhanced,breakable,colback=softolive,colframe=deepolive,
  boxrule=0.7pt,arc=1mm,left=2.3mm,right=2.3mm,top=1.7mm,bottom=1.7mm,
  title={Main result},fonttitle=\bfseries,#1
}
\newtcolorbox{transbox}[1][]{
  enhanced,breakable,colback=softteal,colframe=deepteal,
  boxrule=0.7pt,arc=1mm,left=2.3mm,right=2.3mm,top=1.7mm,bottom=1.7mm,
  title={Transseries form},fonttitle=\bfseries,#1
}
\newtcolorbox{auditbox}[1][]{
  enhanced,breakable,colback=softcoral,colframe=warmgray,
  boxrule=0.6pt,arc=1mm,left=2.3mm,right=2.3mm,top=1.7mm,bottom=1.7mm,
  title={Scope and status},fonttitle=\bfseries,#1
}
\newtcolorbox{algbox}[1][]{
  enhanced,breakable,colback=softblue,colframe=deepblue,
  boxrule=0.6pt,arc=1mm,left=2.3mm,right=2.3mm,top=1.7mm,bottom=1.7mm,
  title={Coefficient engine},fonttitle=\bfseries,#1
}

\newcommand{\Bell}{\mathcal B}
\newcommand{\Tch}{\mathscr T}
\newcommand{\E}{\mathbb E}
\newcommand{\N}{\mathbb N}
\newcommand{\Q}{\mathbb Q}
\newcommand{\C}{\mathbb C}
\newcommand{\R}{\mathbb R}
\newcommand{\ii}{\mathrm i}
\newcommand{\ee}{\mathrm e}
\newcommand{\dd}{\,\mathrm d}
\newcommand{\BigO}{\mathcal O}
\newcommand{\littleo}{\mathrm o}
\newcommand{\coef}[2]{[#1^{#2}]}
\newcommand{\Stwo}[2]{\genfrac\{\}{0pt}{}{#1}{#2}}
\newcommand{\Sone}[2]{\genfrac{[}{]}{0pt}{}{#1}{#2}}
\newcommand{\Gauss}{\mathcal G}
\newcommand{\LambertW}{W}
\newcommand{\Bern}{\beta}
\newcommand{\cA}{\mathsf A}
\newcommand{\cD}{\mathsf D}
\newcommand{\cR}{\mathsf R}
\newcommand{\cH}{\mathsf H}
\newcommand{\cE}{\mathsf E}
\DeclareMathOperator{\Res}{Res}
\DeclareMathOperator{\Log}{Log}
\DeclareMathOperator{\RePart}{Re}
\DeclareMathOperator{\ImPart}{Im}

\title{\Huge\bfseries The Full Asymptotic Transseries of the Bell Numbers\\[3mm]
\Large Finite all-orders saddle coefficients, Lambert--\(W\) geometry,\\
Touchard cumulants, and polynomial--logarithmic re-expansion}
\author{A repository-guided research synthesis with complete coefficient formulae}
\date{September 2026}

\begin{document}
\maketitle

\begin{abstract}
The Bell numbers \(\Bell_n\), defined by
\[
  \exp(\ee^z-1)=\sum_{n\ge 0}\Bell_n\frac{z^n}{n!},
\]
have a classical saddle point at \(r=W_0(n)\).  The leading Moser--Wyman--de
Bruijn approximation is well known, but an actually usable ``all-orders formula''
requires more than unnamed correction polynomials: it must specify a genuine
asymptotic scale, a finite rule for every coefficient, and a remainder theorem.
This article supplies that data.

Write
\[
 \Tch_j(r)=\ee^{-r}(r\partial_r)^j\ee^r
          =\sum_{k=1}^j\Stwo{j}{k}r^k,
 \qquad
 q_j(r)=\frac{\Tch_j(r)}{r(1+r)^{j/2}}.
\]
For every \(m\ge 0\), the \(m\)-th saddle coefficient is given by a finite
partition sum
\[
 c_m(r)=
 \sum_{\substack{k_1,\ldots,k_{2m}\ge0\\
                   \sum_{s=1}^{2m}s k_s=2m}}
 (-1)^{m+K}(2m+2K-1)!!
 \prod_{s=1}^{2m}
 \frac{q_{s+2}(r)^{k_s}}{((s+2)!)^{k_s}k_s!},
 \qquad K=\sum_{s=1}^{2m}k_s,
\]
with \(c_0=1\).  Equivalently, the same coefficients are Gaussian averages of
complete exponential Bell polynomials.  We prove, uniformly to arbitrary fixed
order,
\[
 \Bell_n=
 \frac{n!\exp(\ee^r-1)}{r^n\sqrt{2\pi n(1+r)}}
 \left(\sum_{m=0}^{M-1}\frac{c_m(r)}{n^m}
       +\BigO_M\!\left((r/n)^M\right)\right).
\]
Every \(c_m\) is rational in \(r\), more precisely
\(c_m(r)=P_m(r)/(1+r)^{3m}\) with \(P_m\in\Q[r]\),
\(\deg P_m\le 4m\), and \(c_m(r)=\BigO(r^m)\).

After combining this saddle expansion with the complete Stirling series for
\(n!\), the Bell numbers acquire the natural exponential transseries scale
\(\ee^{-r}=r/n\):
\[
 \log \Bell_n
 \sim n\left(r-1+\frac1r\right)-1-\frac12\log(1+r)
       +\sum_{m\ge1}\cD_m(r)\ee^{-mr},
\]
where every \(\cD_m(r)\) is a bounded rational function given by finite
coefficient operations.  Exponentiating yields a multiplicative transseries with
finite general formulae for its coefficients.  Finally, substituting the
Stirling-cycle-number expansion of \(W_0(n)\) gives the full
polynomial--logarithmic form in \(L=\log n\) and \(\lambda=\log L\), including
the beyond-all-orders sectors \(n^{-m}L^m\Q[\lambda][[L^{-1}]]\).  The same
universal saddle coefficients also govern the Touchard polynomials
\(\Bell_n(x)\) for fixed \(x>0\).
\end{abstract}

\begin{auditbox}
This is a self-contained mathematical article, not a claim of Lean
formalization.  Its principal contribution is an explicit organization of the
standard saddle-point mechanism into finite coefficient formulae and a ranked
transseries.  The first nontrivial correction agrees with the independently
reported 2026 expansion of Hwang, Li, and Zacharovas.  The section on nonprincipal
complex saddles is deliberately labeled \emph{formal/hyperasymptotic}: determining
all Stokes multipliers of the original Cauchy contour is a separate global
steepest-descent problem.
\end{auditbox}

\tableofcontents
\newpage

\section{Introduction and statement of the problem}

\subsection{Bell numbers and the missing all-orders data}

The Bell number \(\Bell_n\) counts set partitions of an \(n\)-element labelled
set.  Its most useful analytic definition is the exponential generating
function
\begin{equation}\label{eq:bell-egf}
  F(z):=\exp(\ee^z-1)
       =\sum_{n\ge0}\Bell_n\frac{z^n}{n!}.
\end{equation}
The coefficient is therefore
\begin{equation}\label{eq:cauchy-basic}
 \frac{\Bell_n}{n!}
 =\frac{1}{2\pi\ii}\int_{|z|=\rho}
   \frac{\exp(\ee^z-1)}{z^{n+1}}\dd z,
 \qquad \rho>0.
\end{equation}
The positive saddle is selected by
\begin{equation}\label{eq:saddle-equation}
   r\ee^r=n,
   \qquad r=W_0(n).
\end{equation}
This gives the classical leading approximation
\begin{equation}\label{eq:classical-leading}
 \Bell_n\sim
 \frac{n!\exp(\ee^r-1)}{r^n\sqrt{2\pi n(1+r)}}
 \sim\frac{\exp\!\left(n(r-1+r^{-1})-1\right)}{\sqrt{1+r}}.
\end{equation}
The two right-hand expressions in \cref{eq:classical-leading} are
asymptotically equivalent; the second uses only the leading factor in
Stirling's formula for \(n!\).

A complete refinement must answer four questions.
\begin{enumerate}[label=\textup{(\roman*)},leftmargin=2.2em]
\item What is the asymptotic scale?
\item How is the coefficient of every term defined by a finite computation?
\item What algebraic structure do those coefficients possess?
\item What remainder remains after an arbitrary fixed truncation?
\end{enumerate}
The linked coefficient-calculus manuscript explicitly identifies these as the
missing obligations behind a schematic Moser--Wyman display.  The purpose here
is to discharge them for the Bell numbers and, at almost no extra cost, for the
Touchard polynomials.

\subsection{Three equivalent notions of ``full expansion''}

There are three useful normal forms, and confusing them obscures the hierarchy
of small terms.

\paragraph{Saddle normal form.}
Keep \(r=W_0(n)\) exact and expand the normalized Cauchy integral in powers of
\(n^{-1}\).  Its coefficients grow like \(r^m\), so the actual term of order
\(m\) is \(\BigO((r/n)^m)\).

\paragraph{Lambert--\(W\) exponential normal form.}
Since \(r/n=\ee^{-r}\), absorb the Stirling series for \(n!\) and rewrite the
answer in powers of the genuine transmonomial \(\ee^{-r}\).  The coefficient
functions are then bounded rational functions of \(r\).

\paragraph{Polynomial--logarithmic normal form.}
Write \(L=\log n\), \(\lambda=\log L\), expand \(r=W_0(n)\) in
\(L^{-1}\Q[\lambda]\), and compose.  The logarithm of \(\Bell_n\) splits into
three ranked sectors:
\[
 nL\,\Q[\lambda][[L^{-1}]],
 \qquad
 \log L+\Q[\lambda][[L^{-1}]],
 \qquad
 \sum_{m\ge1}n^{-m}L^m\Q[\lambda][[L^{-1}]].
\]
The last family is beyond every power of \(L^{-1}\) relative to the first two.
Nonprincipal complex saddles lie farther beyond all these sectors.

\subsection{Road map}

\Cref{sec:exact} develops the exact integral and the Touchard-polynomial
cumulants.  \Cref{sec:allorders} proves the all-orders saddle theorem.
\Cref{sec:coeffcalc} gives four equivalent finite formulae for the coefficients
and proves their rationality.  \Cref{sec:explicit} lists the first coefficients.
\Cref{sec:logtrans} combines the result with Stirling's series and produces the
\(\ee^{-r}\)-transseries.  \Cref{sec:polylog} performs the full
polynomial--logarithmic re-expansion.  \Cref{sec:complex} describes the formal
secondary-saddle sectors, \Cref{sec:touchard} gives the Touchard extension,
\Cref{sec:algorithm} presents exact algorithms, and \Cref{sec:numerics}
quantifies the numerical improvement.

\section{Exact representations and saddle geometry}\label{sec:exact}

\subsection{Combinatorial and probabilistic formulae}

The Stirling numbers of the second kind satisfy
\[
  \Bell_n=\sum_{k=0}^n\Stwo{n}{k}.
\]
Dobinski's formula gives a second exact representation:
\begin{equation}\label{eq:dobinski}
  \Bell_n=\ee^{-1}\sum_{k=0}^{\infty}\frac{k^n}{k!}.
\end{equation}
Indeed,
\[
 \ee^{-1}\sum_{k\ge0}\frac{\ee^{kz}}{k!}
 =\ee^{-1}\exp(\ee^z)=\exp(\ee^z-1),
\]
and comparison with \cref{eq:bell-egf} proves \cref{eq:dobinski}.  Thus
\(\Bell_n\) is also the \(n\)-th raw moment of a Poisson random variable of
mean one.

The saddle method may be applied either to \cref{eq:cauchy-basic} or, after a
local central-limit analysis, to \cref{eq:dobinski}.  The contour route is more
convenient for obtaining every correction coefficient in one uniform formula.

\subsection{The circular integral through the positive saddle}

Set \(z=r\ee^{\ii\theta}\) in \cref{eq:cauchy-basic}, where
\(r=W_0(n)\).  Since \(\dd z=\ii z\dd\theta\),
\begin{equation}\label{eq:circular-integral}
 \frac{\Bell_n}{n!}
 =\frac{\exp(\ee^r-1)}{2\pi r^n}
  \int_{-\pi}^{\pi}\exp\!\bigl(\Phi_r(\theta)\bigr)\dd\theta,
\end{equation}
where
\begin{equation}\label{eq:phase}
 \Phi_r(\theta)
 :=\ee^{r\ee^{\ii\theta}}-\ee^r-\ii n\theta.
\end{equation}
The saddle equation makes \(\Phi_r'(0)=0\).

The complex phase in the original \(z\)-plane is
\begin{equation}\label{eq:zphase}
 \Psi_n(z):=\ee^z-1-n\Log z.
\end{equation}
Its saddle equation \(\Psi_n'(z)=0\) is \(z\ee^z=n\).  Hence all formal saddles
are
\[
 r_k=W_k(n),\qquad k\in\mathbb Z,
\]
but only \(r_0=r>0\) is needed for the ordinary Poincare expansion on the
positive real axis.

\subsection{Touchard polynomials as phase cumulants}

\begin{definition}[Touchard polynomial]\label{def:touchard}
For \(j\ge0\), define
\begin{equation}\label{eq:touchard-def}
 \Tch_j(r):=\ee^{-r}(r\partial_r)^j\ee^r.
\end{equation}
Then \(\Tch_0(r)=1\), and for \(j\ge1\),
\begin{equation}\label{eq:touchard-stirling}
 \Tch_j(r)=\sum_{k=1}^j\Stwo{j}{k}r^k.
\end{equation}
\end{definition}

\begin{proof}
The operator identity
\[
 (r\partial_r)^j=\sum_{k=0}^j\Stwo{j}{k}r^k\partial_r^k
\]
is the standard change from powers to falling factorials.  Applying it to
\(\ee^r\), for which \(\partial_r^k\ee^r=\ee^r\), gives
\cref{eq:touchard-stirling}.
\end{proof}

The first values are
\begin{align*}
 \Tch_1&=r,\\
 \Tch_2&=r+r^2,\\
 \Tch_3&=r+3r^2+r^3,\\
 \Tch_4&=r+7r^2+6r^3+r^4,\\
 \Tch_5&=r+15r^2+25r^3+10r^4+r^5.
\end{align*}
They obey the differential recurrence
\begin{equation}\label{eq:touchard-recurrence}
 \Tch_{j+1}(r)=r\bigl(\Tch_j(r)+\Tch_j'(r)\bigr).
\end{equation}

\begin{lemma}[Exact phase expansion]\label{lem:phase-expansion}
For \(|\theta|\) sufficiently small,
\begin{equation}\label{eq:phase-taylor}
 \Phi_r(\theta)
 =\sum_{j\ge2}\frac{\ii^j\ee^r\Tch_j(r)}{j!}\theta^j.
\end{equation}
In particular, with
\begin{equation}\label{eq:variance-A}
 A:=\ee^r\Tch_2(r)=n(1+r),
\end{equation}
we have \(\Phi_r(\theta)=-A\theta^2/2+\BigO(nr^2|\theta|^3)\).
\end{lemma}

\begin{proof}
Differentiation with respect to \(\theta\) acts on
\(\ee^{r\ee^{\ii\theta}}\) as \(\ii r\partial_r\) at \(\theta=0\).  Therefore
its \(j\)-th derivative there is \(\ii^j\ee^r\Tch_j(r)\).  The linear term is
\(\ii\ee^r\Tch_1(r)\theta=\ii r\ee^r\theta=\ii n\theta\) and cancels the last
term in \cref{eq:phase}.  The quadratic identity follows from
\(\Tch_2=r(1+r)\) and \(r\ee^r=n\).
\end{proof}

\subsection{Standardized cumulants}

Scale
\begin{equation}\label{eq:theta-scale}
  \theta=\frac{u}{\sqrt{A}}
  =\frac{u}{\sqrt{n(1+r)}}.
\end{equation}
For \(s\ge1\), define
\begin{equation}\label{eq:qj-def}
 q_{s+2}(r):=\frac{\Tch_{s+2}(r)}{r(1+r)^{(s+2)/2}},
 \qquad
 a_s(u;r):=\frac{\ii^{s+2}q_{s+2}(r)}{(s+2)!}u^{s+2}.
\end{equation}
Then \cref{eq:phase-taylor} becomes the exact formal identity
\begin{equation}\label{eq:scaled-phase}
 \Phi_r\!\left(\frac{u}{\sqrt A}\right)
 =-\frac{u^2}{2}+\sum_{s\ge1}n^{-s/2}a_s(u;r).
\end{equation}
Since \(\deg \Tch_{s+2}=s+2\),
\begin{equation}\label{eq:q-growth}
  q_{s+2}(r)=\BigO_s(r^{s/2})
  \qquad(r\to\infty).
\end{equation}
Thus the effective small parameter is not merely \(n^{-1/2}\), but
\[
  \varepsilon:=\sqrt{\frac rn}=\ee^{-r/2}.
\]
The Gaussian symmetry removes odd powers of \(\varepsilon\); the final expansion
therefore proceeds in powers of \(\varepsilon^2=r/n=\ee^{-r}\).

\section{The all-orders saddle theorem}\label{sec:allorders}

\subsection{Localization of the Cauchy integral}

The first issue is not formal expansion but localization: the part of the circle
away from \(\theta=0\) must be smaller than every retained algebraic term.

\begin{lemma}[Central and complementary arcs]\label{lem:localization}
There exist absolute constants \(\delta,c_1,c_2>0\) such that, for all
sufficiently large \(n\), with \(r=W_0(n)\),
\begin{align}
 \RePart\Phi_r(\theta)&\le -c_1A\theta^2,
 && |\theta|\le \delta/r,\label{eq:central-decay}\\
 \RePart\Phi_r(\theta)&\le -c_2n/r^2,
 && \delta/r\le |\theta|\le\pi.\label{eq:outer-decay}
\end{align}
Consequently the complementary arc contributes
\begin{equation}\label{eq:outer-small}
 \int_{\delta/r\le|\theta|\le\pi}
       \exp(\Phi_r(\theta))\dd\theta
 =\BigO\!\left(\exp(-c_2n/r^2)\right),
\end{equation}
which is smaller than \((r/n)^M\) for every fixed \(M\).
\end{lemma}

\begin{proof}
For the central bound, \cref{lem:phase-expansion} gives
\[
 \RePart\Phi_r(\theta)
 =-\frac12A\theta^2+\BigO(nr^2|\theta|^3).
\]
Because \(A=n(1+r)\asymp nr\), the ratio of the error to
\(A\theta^2\) is \(\BigO(r|\theta|)\).  Choosing \(\delta\) sufficiently small
proves \cref{eq:central-decay}.

For the complementary bound, discard the oscillating cosine:
\begin{align*}
 \RePart\Phi_r(\theta)
 &=\ee^{r\cos\theta}\cos(r\sin\theta)-\ee^r\\
 &\le \ee^{r\cos\theta}-\ee^r
 =-\ee^r\bigl(1-\ee^{-r(1-\cos\theta)}\bigr).
\end{align*}
On \([-\pi,\pi]\), \(1-\cos\theta\ge C\min(\theta^2,1)\).  If
\(|\theta|\ge\delta/r\), then
\(r(1-\cos\theta)\ge C'\delta^2/r\) until this quantity is bounded away from
zero; in either case
\[
 1-\ee^{-r(1-\cos\theta)}\ge C''/r.
\]
Since \(\ee^r=n/r\), this gives \cref{eq:outer-decay}.  Finally,
\(n/r^2\to\infty\) faster than \(\log(n/r)\), so the exponential in
\cref{eq:outer-small} is smaller than every fixed power of \(r/n\).
\end{proof}

\begin{remark}
The much-studied Hayman admissibility of \(F(z)=\exp(\ee^z-1)\) gives a general
framework for this localization.  The elementary inequalities above are included
so that the particular Bell-number argument does not depend on quoting the full
admissibility machinery.
\end{remark}

\subsection{Gaussian reduction}

Let \(Z\) be standard normal and write
\[
  \E f(Z)=\frac1{\sqrt{2\pi}}
          \int_{-\infty}^{\infty}f(u)\ee^{-u^2/2}\dd u.
\]
After \cref{eq:theta-scale}, the central part of \cref{eq:circular-integral}
becomes
\begin{equation}\label{eq:gaussian-reduction}
 \frac{1}{\sqrt{2\pi A}}
 \int \ee^{-u^2/2}
 \exp\!\left(\sum_{s\ge1}n^{-s/2}a_s(u;r)\right)\dd u,
\end{equation}
up to a contribution smaller than every power of \(r/n\).  Formally replacing
the finite scaled interval by \(\R\) is justified by
\cref{eq:central-decay}; the Gaussian tails are super-algebraically small.

For a precise truncation, choose \(0<\alpha<1/6\) and restrict first to
\(|u|\le U_n:=(n/r)^\alpha\).  On this interval,
\[
 n^{-s/2}a_s(u;r)
 =\BigO_s\!\left((r/n)^{s/2}U_n^{s+2}\right),
\]
and the right side tends to zero for every fixed \(s\).  Taylor's theorem for
the exponential may therefore be applied uniformly to any prescribed order.
Outside \([-U_n,U_n]\), \cref{eq:central-decay} supplies a bound
\(\BigO(\ee^{-cU_n^2})\).  This is the standard Watson-lemma mechanism, here
with coefficients depending slowly on \(r\).

\subsection{All-orders theorem}

\begin{theorem}[Full principal-saddle expansion]\label{thm:all-orders}
Let \(r=W_0(n)\), and define \(a_s\) by \cref{eq:qj-def}.  Put
\begin{equation}\label{eq:cm-gaussian-def}
 c_m(r):=
 \E\!\left[
   \coef{t}{2m}
   \exp\!\left(\sum_{s\ge1}a_s(Z;r)t^s\right)
 \right],
 \qquad m\ge0.
\end{equation}
Then \(c_0(r)=1\), and for every fixed integer \(M\ge1\),
\begin{equation}\label{eq:all-orders-main}
 \boxed{
 \Bell_n=
 \frac{n!\exp(\ee^r-1)}{r^n\sqrt{2\pi n(1+r)}}
 \left(
   \sum_{m=0}^{M-1}\frac{c_m(r)}{n^m}
   +\BigO_M\!\left((r/n)^M\right)
 \right).}
\end{equation}
The constant in the remainder depends on \(M\), but not on \(n\).
\end{theorem}

\begin{proof}
Expand the second exponential in \cref{eq:gaussian-reduction} in the bookkeeping
variable \(t=n^{-1/2}\):
\begin{equation}\label{eq:formal-expansion-Pj}
 \exp\!\left(\sum_{s\ge1}a_s(u;r)t^s\right)
 =\sum_{j\ge0}P_j(u;r)t^j.
\end{equation}
Every monomial contributing to \(P_j\) has the form
\(\prod_s a_s^{k_s}\) with \(\sum_s sk_s=j\).  Its degree in \(u\) is
\[
 \sum_s(s+2)k_s=j+2\sum_s k_s,
\]
which has the same parity as \(j\).  Therefore \(P_j(-u;r)=(-1)^jP_j(u;r)\),
and \(\E P_j(Z;r)=0\) for odd \(j\).  The coefficient of \(n^{-m}\) is
exactly \(c_m(r)=\E P_{2m}(Z;r)\).

By \cref{eq:q-growth}, each monomial in \(P_{2m}\) is
\(\BigO_m(r^m)\) after Gaussian averaging, because
\(\sum_s sk_s=2m\).  Hence
\begin{equation}\label{eq:cm-growth-preliminary}
 c_m(r)=\BigO_m(r^m).
\end{equation}
The uniform Taylor argument preceding the theorem, with enough terms retained,
shows that the first omitted even contribution is
\(\BigO_M(r^M n^{-M})\), while the central Gaussian tail and the complementary
arc are smaller than this by \cref{lem:localization}.  Multiplying by the
prefactor in \cref{eq:circular-integral} yields
\cref{eq:all-orders-main}.
\end{proof}

\begin{masterbox}
The asymptotic scale in \cref{eq:all-orders-main} is genuine:
\[
  1,\quad \frac rn,\quad \left(\frac rn\right)^2,\quad\ldots
  \qquad\left(\frac rn=\ee^{-r}\right).
\]
The displayed coefficient \(c_m/n^m\) has size \(\BigO((r/n)^m)\), and the
remainder after \(M\) terms has the size of the next member of this scale.  No
unnamed correction polynomial is required.
\end{masterbox}

\section{Finite coefficient calculus}\label{sec:coeffcalc}

\subsection{Complete Bell-polynomial formula}

Let \(Y_j\) denote the complete exponential Bell polynomial, characterized by
\begin{equation}\label{eq:complete-bell-def}
 \exp\!\left(\sum_{s\ge1}x_s\frac{t^s}{s!}\right)
 =\sum_{j\ge0}Y_j(x_1,\ldots,x_j)\frac{t^j}{j!}.
\end{equation}
Taking \(x_s=s!a_s(Z;r)\) in \cref{eq:complete-bell-def} gives the first finite
closed rule.

\begin{proposition}[Bell-polynomial coefficient]\label{prop:bellpoly-cm}
For \(m\ge0\),
\begin{equation}\label{eq:cm-bellpoly}
 c_m(r)=\frac1{(2m)!}\E\!
 \left[
  Y_{2m}\bigl(1!a_1(Z;r),2!a_2(Z;r),\ldots,(2m)!a_{2m}(Z;r)\bigr)
 \right].
\end{equation}
Thus \(c_m\) uses only
\(\Tch_3,\Tch_4,\ldots,\Tch_{2m+2}\).
\end{proposition}

\begin{proof}
This is coefficient comparison in \cref{eq:complete-bell-def}; terms with
\(s>2m\) cannot affect \(\coef{t}{2m}\).
\end{proof}

The Gaussian average can itself be written as a differential operator:
\begin{equation}\label{eq:gaussian-heat}
 \E P(Z)=\left.\exp\!\left(\frac12\partial_u^2\right)P(u)\right|_{u=0}
\end{equation}
for every polynomial \(P\).  Consequently \cref{eq:cm-bellpoly} is a finite
composition of two standard coefficient-calculus operations: a complete Bell
polynomial followed by the heat operator at the origin.

\subsection{Explicit partition sum}

The Bell-polynomial formula becomes completely elementary after expanding the
exponential and using
\(\E Z^{2j}=(2j-1)!!\), \(\E Z^{2j+1}=0\).

\begin{theorem}[Finite partition formula]\label{thm:partition-cm}
For \(m\ge1\),
\begin{equation}\label{eq:cm-partition}
 \boxed{
 c_m(r)=
 \sum_{\substack{k_1,\ldots,k_{2m}\ge0\\
                   \sum_{s=1}^{2m}s k_s=2m}}
 (-1)^{m+K}(2m+2K-1)!!
 \prod_{s=1}^{2m}
 \frac{q_{s+2}(r)^{k_s}}{((s+2)!)^{k_s}k_s!},}
 \qquad K:=\sum_{s=1}^{2m}k_s.
\end{equation}
We use the convention \((-1)!!=1\).
\end{theorem}

\begin{proof}
The coefficient of \(t^{2m}\) in
\(\exp(\sum_sa_st^s)\) is
\[
 \sum_{\sum sk_s=2m}\prod_s\frac{a_s^{k_s}}{k_s!}.
\]
The power of \(Z\) in such a term is
\(\sum(s+2)k_s=2m+2K\), and the power of \(\ii\) is the same.  Hence
\[
 \ii^{2m+2K}=(-1)^{m+K},
 \qquad
 \E Z^{2m+2K}=(2m+2K-1)!!.
\]
Substitution of \cref{eq:qj-def} proves \cref{eq:cm-partition}.
\end{proof}

\begin{remark}[Finiteness]
The sum is indexed by partitions of \(2m\): the vector
\((k_1,\ldots,k_{2m})\) records \(k_s\) parts of size \(s\).  Therefore the
number of summands is the partition number \(p(2m)\), not an infinite or
asymptotic count.  This is the requested finite coefficient-extraction rule.
\end{remark}

\subsection{A dynamic recurrence}

Direct enumeration of partitions is conceptually transparent.  For symbolic
computation, a recurrence for the intermediate polynomials is often faster.
Differentiating \cref{eq:formal-expansion-Pj} with respect to \(t\) gives
\begin{equation}\label{eq:Pj-recurrence}
 P_0(u;r)=1,
 \qquad
 P_j(u;r)=\frac1j\sum_{s=1}^{j}s\,a_s(u;r)P_{j-s}(u;r)
 \quad(j\ge1).
\end{equation}
Then
\begin{equation}\label{eq:cm-from-P}
 c_m(r)=\E P_{2m}(Z;r).
\end{equation}
Because parity is preserved by \cref{eq:Pj-recurrence}, odd \(P_j\) need not be
Gaussian-averaged at all.

\subsection{Rationality, denominators, and degrees}

\begin{theorem}[Algebraic structure of the coefficients]\label{thm:cm-structure}
For every \(m\ge0\), there is a polynomial \(P_m(r)\in\Q[r]\) such that
\begin{equation}\label{eq:cm-rational-structure}
 c_m(r)=\frac{P_m(r)}{(1+r)^{3m}},
 \qquad
 \deg P_m\le4m.
\end{equation}
Moreover
\begin{equation}\label{eq:cm-growth}
 c_m(r)=\BigO_m(r^m)\qquad(r\to\infty).
\end{equation}
\end{theorem}

\begin{proof}
Fix a vector in \cref{eq:cm-partition}.  Since
\(\sum sk_s=2m\),
\begin{equation}\label{eq:denom-exponent}
 \prod_s(1+r)^{(s+2)k_s/2}
 =(1+r)^{m+K}.
\end{equation}
Also \(K\le2m\), because every part has size at least one.  Bringing the term to
the common denominator \((1+r)^{3m}\) multiplies its numerator by
\((1+r)^{2m-K}\).

Now \(\Tch_{s+2}(r)/r\) is a polynomial of degree \(s+1\).  The product of
these polynomials has degree
\[
 \sum_s(s+1)k_s=2m+K.
\]
After multiplication by \((1+r)^{2m-K}\), the degree is at most \(4m\).
This proves \cref{eq:cm-rational-structure}.  Finally,
\(q_{s+2}(r)=\BigO(r^{s/2})\), so every summand in
\cref{eq:cm-partition} is \(\BigO(r^{\sum sk_s/2})=\BigO(r^m)\), proving
\cref{eq:cm-growth}.
\end{proof}

\begin{corollary}[Asymptotic sequence]\label{cor:genuine-scale}
For fixed \(m\),
\[
 \frac{c_{m+1}(r)n^{-m-1}}{c_m(r)n^{-m}}
 =\BigO_m(r/n)=\BigO_m(\ee^{-r})
\]
whenever the denominator is nonzero.  Thus the sequence of term sizes is
strictly asymptotic.
\end{corollary}

\subsection{The limiting Stirling identity}

The highest-degree part of \(c_m\) has a classical interpretation.  Define
\begin{equation}\label{eq:gamma-m-def}
 \gamma_m:=\lim_{r\to\infty}\frac{c_m(r)}{r^m}.
\end{equation}
Because \(q_{s+2}(r)/r^{s/2}\to1\), \cref{eq:cm-gaussian-def} yields
\begin{equation}\label{eq:gamma-gaussian}
 \gamma_m=
 \E\!\left[
  \coef{t}{2m}
  \exp\!\left(
    \sum_{s\ge1}\frac{\ii^{s+2}Z^{s+2}}{(s+2)!}t^s
  \right)
 \right].
\end{equation}
Comparison with the saddle derivation of Stirling's formula gives the formal
identity
\begin{equation}\label{eq:gamma-stirling-inverse}
 \boxed{
 \sum_{m\ge0}\gamma_m x^m
 =\exp\!\left(
   -\sum_{j\ge1}\frac{\Bern_{2j}}{2j(2j-1)}x^{2j-1}
 \right),}
\end{equation}
where
\begin{equation}\label{eq:bernoulli-convention}
 \frac{t}{\ee^t-1}=\sum_{j\ge0}\Bern_j\frac{t^j}{j!},
 \qquad \Bern_1=-\frac12.
\end{equation}
This is the Gaussian coefficient identity that appears when Stirling's formula
is derived on two equivalent saddle contours.

The first values are
\begin{equation}\label{eq:gamma-values}
 \gamma_0=1,
 \quad \gamma_1=-\frac1{12},
 \quad \gamma_2=\frac1{288},
 \quad \gamma_3=\frac{139}{51840},
 \quad \gamma_4=-\frac{571}{2488320}.
\end{equation}
Equation \cref{eq:gamma-stirling-inverse} supplies a useful independent check on
any symbolic implementation of \cref{eq:cm-partition}.

\section{Explicit principal-saddle coefficients}\label{sec:explicit}

\subsection{The first correction by hand}

Only \(a_1\) and \(a_2\) can contribute to \(c_1\):
\[
 c_1=\E\!\left(a_2(Z;r)+\frac12a_1(Z;r)^2\right).
\]
Since
\[
 a_1=-\frac{\ii q_3}{6}u^3,
 \qquad
 a_2=\frac{q_4}{24}u^4,
\]
and \(\E Z^4=3\), \(\E Z^6=15\),
\begin{equation}\label{eq:c1-q}
 c_1(r)=\frac18q_4(r)-\frac5{24}q_3(r)^2.
\end{equation}
Substitution of \(\Tch_3\) and \(\Tch_4\) gives
\begin{equation}\label{eq:c1-explicit}
 \boxed{
 c_1(r)=
 -\frac{2r^4+9r^3+16r^2+6r+2}{24(1+r)^3}.}
\end{equation}
The sign is worth noting: for large \(r\), \(c_1(r)\sim-r/12\).

\subsection{The second correction}

At order \(t^4\), the partitions of four give
\begin{align}
 c_2={}&-\frac1{48}q_6
 +\frac7{48}q_3q_5
 +\frac{35}{384}q_4^2
 -\frac{35}{64}q_3^2q_4
 +\frac{385}{1152}q_3^4.\label{eq:c2-q}
\end{align}
After rational simplification,
\begin{equation}\label{eq:c2-explicit}
 \boxed{
 c_2(r)=
 \frac{
 4r^8-12r^7-143r^6-384r^5-588r^4-636r^3
 +100r^2+24r+4}
 {1152(1+r)^6}.}
\end{equation}

\subsection{The third correction}

The finite rule \cref{eq:cm-partition} gives
\begin{equation}\label{eq:c3-explicit}
\boxed{
\begin{aligned}
 c_3(r)=\frac{1}{414720(1+r)^9}\bigl(&
 1112r^{12}+12708r^{11}+56082r^{10}+118683r^9\\
 &+97176r^8-74394r^7-259046r^6-456444r^5\\
 &-635424r^4+115308r^3+39432r^2+10008r+1112
 \bigr).
\end{aligned}}
\end{equation}
The leading coefficients of \(c_1,c_2,c_3\) agree with
\(\gamma_1,\gamma_2,\gamma_3\) in \cref{eq:gamma-values}.

\subsection{How to read the first terms}

Substituting \cref{eq:c1-explicit,eq:c2-explicit,eq:c3-explicit} into
\cref{eq:all-orders-main} gives
\begin{align}\label{eq:bell-first-cs}
 \Bell_n={}&
 \frac{n!\exp(\ee^r-1)}{r^n\sqrt{2\pi n(1+r)}}
 \left(
  1+\frac{c_1(r)}n+\frac{c_2(r)}{n^2}
    +\frac{c_3(r)}{n^3}
    +\BigO\!\left(\frac{r^4}{n^4}\right)
 \right).
\end{align}
This is already a complete Poincare expansion, but it hides the most natural
transmonomial because \(c_m(r)\) grows like \(r^m\).  The next section absorbs
Stirling's series and rewrites every term as a bounded coefficient times
\(\ee^{-mr}\).

\section{Logarithmic and exponential transseries}\label{sec:logtrans}

\subsection{Combining the saddle series with Stirling's series}

Stirling's complete logarithmic expansion is
\begin{equation}\label{eq:stirling-log}
 \log n!
 \sim n\log n-n+\frac12\log(2\pi n)
 +\sum_{j\ge1}\frac{\Bern_{2j}}{2j(2j-1)n^{2j-1}}.
\end{equation}
Define
\begin{equation}\label{eq:sigma-m}
 \sigma_m:=
 \begin{cases}
 \displaystyle\frac{\Bern_{m+1}}{m(m+1)},&m\text{ odd},\\[2mm]
 0,&m\text{ even}.
 \end{cases}
\end{equation}
Thus the last sum in \cref{eq:stirling-log} is
\(\sum_{m\ge1}\sigma_mn^{-m}\).

Let
\begin{equation}\label{eq:ell-def}
 C(x;r):=1+\sum_{m\ge1}c_m(r)x^m,
 \qquad
 \log C(x;r)=\sum_{m\ge1}\ell_m(r)x^m.
\end{equation}
Every \(\ell_m\) is finite and can be obtained in either of the forms
\begin{align}
 \ell_m(r)
 &=\sum_{\substack{k_1,\ldots,k_m\ge0\\
                    \sum_{j=1}^mjk_j=m}}
   (-1)^{K+1}(K-1)!
   \prod_{j=1}^m\frac{c_j(r)^{k_j}}{k_j!},
 \qquad K=\sum_jk_j,
 \label{eq:ell-partition}\\
 \ell_1&=c_1,
 \qquad
 \ell_m=c_m-\frac1m\sum_{j=1}^{m-1}j\ell_jc_{m-j}
 \quad(m\ge2).
 \label{eq:ell-recurrence}
\end{align}
The recurrence follows from \(C'=C(\log C)'\).

\begin{definition}[Logarithmic saddle coefficients]\label{def:d-m}
Set
\begin{equation}\label{eq:d-m-def}
 d_m(r):=\sigma_m+\ell_m(r).
\end{equation}
\end{definition}

\begin{theorem}[Lambert--\(W\) logarithmic form]\label{thm:log-W-form}
For every fixed \(M\ge1\), with \(r=W_0(n)\),
\begin{equation}\label{eq:log-W-truncated}
 \boxed{
 \log\Bell_n
 =n\left(r-1+\frac1r\right)-1-\frac12\log(1+r)
 +\sum_{m=1}^{M-1}\frac{d_m(r)}{n^m}
 +\BigO_M\!\left((r/n)^M\right).}
\end{equation}
Equivalently,
\begin{equation}\label{eq:bell-exp-d}
 \Bell_n=
 \frac{\exp\!\left(n(r-1+r^{-1})-1\right)}{\sqrt{1+r}}
 \exp\!\left(
    \sum_{m=1}^{M-1}\frac{d_m(r)}{n^m}
    +\BigO_M((r/n)^M)
 \right).
\end{equation}
Each \(d_m\) has the form
\begin{equation}\label{eq:d-structure}
 d_m(r)=\frac{Q_m(r)}{(1+r)^{3m}},
 \qquad Q_m\in\Q[r],\qquad \deg Q_m\le4m,
\end{equation}
and \(d_m(r)=\BigO_m(r^m)\).
\end{theorem}

\begin{proof}
Insert \cref{eq:stirling-log} and \cref{eq:ell-def} into the logarithm of
\cref{eq:all-orders-main}.  The identities
\[
 \log n-\log r=r,
 \qquad
 \ee^r=\frac nr
\]
follow from \(r\ee^r=n\) and produce the first line of
\cref{eq:log-W-truncated}.  The structure statement follows from
\cref{thm:cm-structure}, the finite logarithm formula
\cref{eq:ell-partition}, and a common denominator \((1+r)^{3m}\).
\end{proof}

\subsection{The genuine \texorpdfstring{\(\ee^{-r}\)}{exp(-r)}-transseries}

Define the bounded rational functions
\begin{equation}\label{eq:D-m-def}
 \cD_m(r):=\frac{d_m(r)}{r^m}.
\end{equation}
Since \(n=r\ee^r\),
\[
 \frac{d_m(r)}{n^m}=\cD_m(r)\ee^{-mr}.
\]
Therefore \cref{eq:log-W-truncated} becomes the genuine exponential
transseries.

\begin{transbox}
For every fixed \(M\),
\begin{equation}\label{eq:log-e-r-transseries}
 \boxed{
 \begin{aligned}
 \log\Bell_n={}&n\left(r-1+\frac1r\right)-1-\frac12\log(1+r)\\
 &+\sum_{m=1}^{M-1}\cD_m(r)\ee^{-mr}
 +\BigO_M(\ee^{-Mr}),
 \qquad r=W_0(n).
 \end{aligned}}
\end{equation}
Every \(\cD_m\) is rational and bounded as \(r\to\infty\).  Thus
\(1,\ee^{-r},\ee^{-2r},\ldots\) is an honest asymptotic scale with bounded
coefficient functions.
\end{transbox}

The first logarithmic coefficients simplify markedly:
\begin{align}
 d_1(r)
 &=-\frac{r^2(2r^2+7r+10)}{24(1+r)^3},
 \label{eq:d1}\\
 d_2(r)
 &=-\frac{r^3(2r^4+12r^3+29r^2+40r+36)}{48(1+r)^6},
 \label{eq:d2}\\
 d_3(r)
 &=\frac{r^4}{5760(1+r)^9}
 \bigl(16r^8+164r^7+616r^6+818r^5-1096r^4
 \notag\\[-1mm]
 &\hspace{38mm}{}-5793r^3-10032r^2-12208r-12048\bigr).
 \label{eq:d3}
\end{align}
Consequently
\begin{align}
 \cD_1(r)
 &=-\frac{r(2r^2+7r+10)}{24(1+r)^3},
 \label{eq:D1}\\
 \cD_2(r)
 &=-\frac{r(2r^4+12r^3+29r^2+40r+36)}{48(1+r)^6},
 \label{eq:D2}\\
 \cD_3(r)
 &=\frac{r}{5760(1+r)^9}
 \bigl(16r^8+164r^7+616r^6+818r^5-1096r^4
 \notag\\[-1mm]
 &\hspace{38mm}{}-5793r^3-10032r^2-12208r-12048\bigr).
 \label{eq:D3}
\end{align}
Their large-\(r\) expansions begin
\begin{align}
 \cD_1(r)&=-\frac1{12}-\frac1{24r}-\frac1{24r^2}
             +\frac1{3r^3}-\frac5{6r^4}+\cdots,
 \label{eq:D1-invr}\\
 \cD_2(r)&=-\frac1{24r}+\frac1{48r^3}-\frac1{8r^4}+\cdots,
 \label{eq:D2-invr}\\
 \cD_3(r)&=\frac1{360}+\frac1{288r}-\frac7{288r^2}
             +\frac7{2880r^3}+\cdots.
 \label{eq:D3-invr}
\end{align}
Because the only finite singularities of these displayed rational functions are
at \(r=0\) and \(r=-1\), their inverse-\(r\) expansions converge for
sufficiently large \(|r|\), not merely formally.

\subsection{The first correction and comparison with the literature}

Keeping only \(d_1/n\) in \cref{eq:bell-exp-d} gives
\begin{equation}\label{eq:first-literature-correction}
 \boxed{
 \Bell_n=
 \frac{\exp\!\left(n(r-1+r^{-1})-1\right)}{\sqrt{1+r}}
 \left(
  1-
  \frac{r^2(2r^2+7r+10)}{24(1+r)^3n}
  +\BigO\!\left(\frac{r^2}{n^2}\right)
 \right).}
\end{equation}
This is exactly the first correction reported in the 2026 expansion of Hwang,
Li, and Zacharovas, with their notation \(\omega_n=W(n)\).  Here it is not an
isolated correction: it is the first member of the finite all-orders mechanism
\cref{eq:cm-partition,eq:ell-partition,eq:d-m-def}.

\subsection{Multiplicative exponential coefficients}

Define \(\cA_0(r)=1\) and
\begin{equation}\label{eq:A-generating}
 \exp\!\left(\sum_{j\ge1}\cD_j(r)y^j\right)
 =\sum_{m\ge0}\cA_m(r)y^m.
\end{equation}
Then
\begin{align}
 \cA_m(r)
 &=\sum_{\substack{k_1,\ldots,k_m\ge0\\
                    \sum_{j=1}^mjk_j=m}}
   \prod_{j=1}^m\frac{\cD_j(r)^{k_j}}{k_j!},
 \label{eq:A-partition}\\
 m\cA_m(r)
 &=\sum_{j=1}^m j\cD_j(r)\cA_{m-j}(r).
 \label{eq:A-recurrence}
\end{align}
Exponentiating \cref{eq:log-e-r-transseries} gives
\begin{equation}\label{eq:multiplicative-e-r}
 \boxed{
 \Bell_n=
 \frac{\exp\!\left(n(r-1+r^{-1})-1\right)}{\sqrt{1+r}}
 \left(
   \sum_{m=0}^{M-1}\cA_m(r)\ee^{-mr}
   +\BigO_M(\ee^{-Mr})
 \right).}
\end{equation}
The first three bounded coefficient functions are
\begin{align}
 \cA_1(r)
 &=-\frac{r(2r^2+7r+10)}{24(1+r)^3},
 \label{eq:A1}\\
 \cA_2(r)
 &=\frac{r(4r^5-20r^4-199r^3-556r^2-860r-864)}
 {1152(1+r)^6},
 \label{eq:A2}\\
 \cA_3(r)
 &=\frac{r}{414720(1+r)^9}
 \bigl(1112r^8+12828r^7+55962r^6+111301r^5
 \notag\\[-1mm]
 &\hspace{27mm}{}+55818r^4-196476r^3-492584r^2
 -749376r-867456\bigr).
 \label{eq:A3}
\end{align}

\subsection{Bernoulli boundary values}

Taking the logarithm of \cref{eq:gamma-stirling-inverse} identifies the limits
of the bounded logarithmic coefficients.

\begin{proposition}[Limits at the transseries boundary]\label{prop:D-limits}
For \(m\ge1\),
\begin{equation}\label{eq:D-limit}
 \lim_{r\to\infty}\cD_m(r)=
 \begin{cases}
 \displaystyle-\frac{\Bern_{m+1}}{m(m+1)},&m\text{ odd},\\[2mm]
 0,&m\text{ even}.
 \end{cases}
\end{equation}
Moreover
\begin{equation}\label{eq:A-limit}
 \lim_{r\to\infty}\cA_m(r)=\gamma_m.
\end{equation}
\end{proposition}

\begin{proof}
By \cref{eq:ell-def}, the coefficient of \(x^m\) in
\(\log\sum_jc_j(r)x^j\), after the scaling \(x\mapsto x/r\), tends to the
coefficient of \(x^m\) in \(\log\sum_j\gamma_jx^j\).  Equation
\cref{eq:gamma-stirling-inverse} gives \cref{eq:D-limit}; the fixed Stirling
term \(\sigma_m/r^m\) tends to zero.  Exponentiating proves
\cref{eq:A-limit}.
\end{proof}

\section{Polynomial--logarithmic re-expansion}\label{sec:polylog}

\subsection{The all-orders Lambert--\texorpdfstring{\(W\)}{W} coefficients}

Set
\begin{equation}\label{eq:L-lambda}
  L:=\log n,
  \qquad \lambda:=\log L.
\end{equation}
The principal Lambert function has the Stirling-cycle-number expansion
\begin{equation}\label{eq:W-polylog}
 r=W_0(n)
 \sim L-\lambda+\sum_{N\ge1}\frac{\cR_N(\lambda)}{L^N},
\end{equation}
where
\begin{equation}\label{eq:R-N-general}
 \boxed{
 \cR_N(\lambda)=
 \sum_{m=1}^{N}
 \frac{(-1)^{N-m}}{m!}
 \Sone{N}{N-m+1}\lambda^m.}
\end{equation}
Here \(\Sone{N}{j}\) is the unsigned Stirling cycle number.  The first
polynomials are
\begin{align}
 \cR_1&=\lambda,\notag\\
 \cR_2&=\frac{\lambda(\lambda-2)}2,\notag\\
 \cR_3&=\frac{\lambda(2\lambda^2-9\lambda+6)}6,\notag\\
 \cR_4&=\frac{\lambda(3\lambda^3-22\lambda^2+36\lambda-12)}{12},
 \label{eq:R-first}\\
 \cR_5&=\frac{\lambda(12\lambda^4-125\lambda^3+350\lambda^2
                         -300\lambda+60)}{60}.\notag
\end{align}
Formula \cref{eq:R-N-general} is the exact coefficient bridge between the
combinatorial first-kind Stirling calculus and the logarithmic transseries
calculus.

\subsection{The extensive sector: coefficients multiplying \texorpdfstring{\(n\)}{n}}

Put \(z=L^{-1}\) and define
\begin{equation}\label{eq:Dformal}
 \Delta(z;\lambda)
 :=1-\lambda z+\sum_{j\ge1}\cR_j(\lambda)z^{j+1}.
\end{equation}
Then \(r=z^{-1}\Delta(z;\lambda)\).  Define
\begin{equation}\label{eq:H-general}
 \boxed{
 \cH_N(\lambda):=
 \cR_N(\lambda)+\coef{z}{N-1}\Delta(z;\lambda)^{-1}.}
\end{equation}
Only \(\cR_1,\ldots,\cR_N\) are needed to calculate \(\cH_N\).  Since
\(1/r=z\Delta^{-1}\),
\begin{equation}\label{eq:F-polylog-general}
 r-1+\frac1r
 \sim L-\lambda-1+
       \sum_{N\ge1}\frac{\cH_N(\lambda)}{L^N}.
\end{equation}
The first coefficients are
\begin{align}
 \cH_1&=\lambda+1,\notag\\
 \cH_2&=\frac{\lambda^2}{2},\notag\\
 \cH_3&=\frac{\lambda^2(2\lambda-3)}6,\notag\\
 \cH_4&=\frac{\lambda^2(3\lambda^2-10\lambda+6)}{12},
 \label{eq:H-first}\\
 \cH_5&=\frac{\lambda^2(12\lambda^3-65\lambda^2+90\lambda-30)}{60},\notag\\
 \cH_6&=\frac{\lambda^2(20\lambda^4-154\lambda^3+355\lambda^2
                              -280\lambda+60)}{120}.\notag
\end{align}
Therefore the classical de Bruijn expansion is refined to arbitrary order by
\begin{equation}\label{eq:normalized-log-debruijn}
 \frac{\log\Bell_n}{n}
 \sim L-\lambda-1+
 \sum_{N\ge1}\frac{\cH_N(\lambda)}{L^N},
\end{equation}
where all omitted sectors are \(\BigO(\log L/n)\), hence beyond every power of
\(L^{-1}\).

Written out through \(L^{-6}\),
\begin{align}\label{eq:debruijn-many}
 \frac{\log\Bell_n}{n}
 \sim{}&L-\lambda-1+\frac{\lambda+1}{L}
 +\frac{\lambda^2}{2L^2}
 +\frac{\lambda^2(2\lambda-3)}{6L^3}\\
 &+\frac{\lambda^2(3\lambda^2-10\lambda+6)}{12L^4}
 +\frac{\lambda^2(12\lambda^3-65\lambda^2+90\lambda-30)}{60L^5}
 \notag\\
 &+\frac{\lambda^2(20\lambda^4-154\lambda^3+355\lambda^2
                              -280\lambda+60)}{120L^6}+\cdots.\notag
\end{align}

\subsection{The edge sector}

The factor \((1+r)^{-1/2}\) lives on a lower transseries level than the terms
multiplied by \(n\).  Define
\begin{equation}\label{eq:Sigmaformal}
 \Sigma(z;\lambda)
 :=z(1+r)=\Delta(z;\lambda)+z
 =1+(1-\lambda)z+\sum_{j\ge1}\cR_j(\lambda)z^{j+1},
\end{equation}
and
\begin{equation}\label{eq:E-general}
 \boxed{
 \cE_N(\lambda):=-\frac12\coef{z}{N}\log\Sigma(z;\lambda).}
\end{equation}
Then
\begin{equation}\label{eq:edge-general}
 -1-\frac12\log(1+r)
 \sim -1-\frac12\log L+
       \sum_{N\ge1}\frac{\cE_N(\lambda)}{L^N}.
\end{equation}
The first values are
\begin{align}
 \cE_1&=\frac{\lambda-1}{2},\notag\\
 \cE_2&=\frac{\lambda^2-4\lambda+1}{4},\notag\\
 \cE_3&=\frac{2\lambda^3-15\lambda^2+18\lambda-2}{12},
 \label{eq:E-first}\\
 \cE_4&=\frac{3\lambda^4-34\lambda^3+84\lambda^2-48\lambda+3}{24},\notag\\
 \cE_5&=\frac{12\lambda^5-185\lambda^4+730\lambda^3-900\lambda^2
                    +300\lambda-12}{120}.\notag
\end{align}
As with \(\ell_m\), an explicit partition formula follows by expanding
\(\log(1+u)\); \cref{eq:E-general} is usually the shortest exact definition.

\subsection{All exponentially small sectors}

For each \(m\ge1\), define polynomial--logarithmic coefficient polynomials
\(p_{m,j}(\lambda)\) by
\begin{equation}\label{eq:p-mj}
 \boxed{
 p_{m,j}(\lambda):=
 \coef{z}{j}
 \left\{
 z^m d_m\!\left(z^{-1}\Delta(z;\lambda)\right)
 \right\}.}
\end{equation}
Because \(d_m(r)=\BigO(r^m)\), the expression in braces is a formal power
series in \(z\), so \cref{eq:p-mj} is well-defined.  It gives
\begin{equation}\label{eq:dm-polylog-sector}
 \frac{d_m(r)}{n^m}
 \sim \frac{L^m}{n^m}
       \sum_{j\ge0}\frac{p_{m,j}(\lambda)}{L^j}.
\end{equation}
Every \(p_{m,j}\) is obtained by a finite computation from Stirling numbers of
both kinds, Gaussian moments, and ordinary coefficient extraction.

For illustration, the first sector begins
\begin{align}\label{eq:first-exp-sector}
 \frac{d_1(r)}n
 \sim \frac{L}{n}\biggl(&
 -\frac1{12}
 +\frac{2\lambda-1}{24L}
 -\frac{2\lambda+1}{24L^2}
 -\frac{\lambda^2-\lambda-8}{24L^3}\\
 &-\frac{2\lambda^3-6\lambda^2-45\lambda+60}{72L^4}
 +\cdots\biggr).\notag
\end{align}
The next two begin
\begin{align}
 \frac{d_2(r)}{n^2}
 &\sim \frac{L^2}{n^2}\left(
 -\frac1{24L}+\frac{\lambda}{24L^2}
 -\frac{2\lambda-1}{48L^3}
 -\frac{\lambda^2-3\lambda+6}{48L^4}+\cdots\right),
 \label{eq:second-exp-sector}\\
 \frac{d_3(r)}{n^3}
 &\sim \frac{L^3}{n^3}\left(
 \frac1{360}-\frac{12\lambda-5}{1440L}
 +\frac{12\lambda^2+2\lambda-35}{1440L^2}+\cdots\right).
 \label{eq:third-exp-sector}
\end{align}

\subsection{The full polynomial--logarithmic transseries}

Combining \cref{eq:F-polylog-general,eq:edge-general,eq:dm-polylog-sector}
gives the promised all-level expansion.

\begin{theorem}[Full logarithmic transseries in \(L\) and \(\lambda\)]
\label{thm:full-polylog}
With \(L=\log n\), \(\lambda=\log L\),
\begin{equation}\label{eq:full-polylog-box}
\boxed{
\begin{aligned}
 \log\Bell_n\sim{}&
 n\left(
   L-\lambda-1+
   \sum_{N\ge1}\frac{\cH_N(\lambda)}{L^N}
 \right)
 -1-\frac12\log L
 +\sum_{N\ge1}\frac{\cE_N(\lambda)}{L^N}\\
 &+\sum_{m\ge1}\frac{L^m}{n^m}
   \sum_{j\ge0}\frac{p_{m,j}(\lambda)}{L^j}.
\end{aligned}}
\end{equation}
The coefficient families are given explicitly by
\cref{eq:R-N-general,eq:H-general,eq:E-general,eq:p-mj}, while the
\(d_m\) entering \cref{eq:p-mj} are defined by the finite formulae
\cref{eq:cm-partition,eq:ell-partition,eq:d-m-def}.
\end{theorem}

\begin{proof}
The first two lines are formal composition of \cref{eq:W-polylog} with the
rational functions in \cref{eq:log-W-truncated}.  Because the inner series has
leading monomial \(L\), inversion, logarithm, and rational composition are
valid in the polynomial--logarithmic transseries ring.  The coefficient
extraction \cref{eq:p-mj} handles each exponentially small sector separately.
The hierarchy is well ordered by first comparing the power of \(n^{-1}\), then
the power of \(L^{-1}\) within that sector.
\end{proof}

\begin{remark}[Exponentiating the full logarithmic transseries]
Equation \cref{eq:full-polylog-box} is usually the cleanest form because
\(\Bell_n\) is enormous.  The Bell number itself is obtained simply as
\(\Bell_n\sim\exp(\text{right-hand side})\).  Alternatively, exponentiate the
last line sector by sector using complete Bell polynomials; this reproduces the
bounded multiplicative coefficients \(\cA_m\) of
\cref{eq:multiplicative-e-r} after the substitution \(r=W_0(n)\).
\end{remark}

\section{Nonprincipal saddles and hyperasymptotic sectors}\label{sec:complex}

\subsection{Universal local germs}

The phase \(\Psi_n(z)=\ee^z-1-n\Log z\) has saddles
\(r_k=W_k(n)\).  Around any nondegenerate saddle, the same local calculation
applies after analytic continuation.  Thus the formal local germ is
\begin{equation}\label{eq:complex-local-germ}
 \mathfrak S_k(n)\sim
 \frac{n!\exp(\ee^{r_k}-1)}{r_k^n\sqrt{2\pi n(1+r_k)}}
 \sum_{m\ge0}\frac{c_m(r_k)}{n^m},
\end{equation}
with a consistent branch of the square root chosen along the steepest-descent
direction.  The coefficient functions \(c_m\) need no modification: they are
rational and hence analytically continue away from \(r_k=-1\).

\subsection{Relative actions}

Let
\begin{equation}\label{eq:Delta-k}
 \Delta_k(n):=\Psi_n(r_k)-\Psi_n(r_0).
\end{equation}
For fixed \(k\), put \(y=2\pi\ii k\).  The relation
\(r_k+\Log r_k=r+\log r+y\) may be written
\begin{equation}\label{eq:delta-eq}
 \delta_k+\log\!\left(1+\frac{\delta_k}{r}\right)=y,
 \qquad \delta_k:=r_k-r.
\end{equation}
Formal reversion gives
\begin{equation}\label{eq:delta-expansion}
 \delta_k
 =y-\frac yr+\frac{y(y+2)}{2r^2}
 +\frac{y(-2y^2-9y-6)}{6r^3}+\BigO_k(r^{-4}).
\end{equation}
Using \(\ee^r=n/r\), one obtains
\begin{equation}\label{eq:action-expansion}
 \frac{\Delta_k(n)}n
 =-\frac yr+\frac{y^2}{2r^2}
 -\frac{y^2/2+y^3/3}{r^3}+\BigO_k(r^{-4}).
\end{equation}
Hence
\begin{equation}\label{eq:real-action-gap}
 \RePart\Delta_k(n)
 =-\frac{2\pi^2k^2n}{r^2}
   \left(1-\frac1r+\BigO_k(r^{-2})\right).
\end{equation}
The adjacent complex saddles are therefore suppressed on the scale
\begin{equation}\label{eq:complex-scale}
 \exp\!\left(-\frac{2\pi^2k^2n}{r^2}
              (1+\BigO_k(r^{-1}))\right),
\end{equation}
which is smaller than \(\ee^{-Mr}\) for every fixed \(M\).

\begin{auditbox}[title={Why no unconditional global saddle sum is asserted}]
A formal expression
\[
  \sum_{k\in\mathbb Z}\sigma_k(n)\mathfrak S_k(n)
\]
requires the intersection numbers or Stokes multipliers \(\sigma_k(n)\) of the
original Cauchy contour with the steepest-descent thimbles.  The local coefficient
calculus determines \(\mathfrak S_k\), but not those global multipliers.  For the
ordinary positive-real Poincare expansion, \(k=0\) alone suffices and the
remainder theorem of \cref{thm:all-orders} is unconditional.  The scales
\cref{eq:complex-scale} describe the natural hyperasymptotic remainder and the
location of possible exponentially tiny oscillations, not an unproved exact
transseries decomposition.
\end{auditbox}

\section{Touchard polynomials and shifted Bell numbers}\label{sec:touchard}

\subsection{Universality for \texorpdfstring{\(\Bell_n(x)\)}{Bell polynomials}}

Define the Touchard polynomial
\begin{equation}\label{eq:touchard-egf-x}
 \exp\!\bigl(x(\ee^z-1)\bigr)
 =\sum_{n\ge0}\Bell_n(x)\frac{z^n}{n!},
 \qquad
 \Bell_n(x)=\sum_{k=0}^n\Stwo{n}{k}x^k.
\end{equation}
For fixed \(x>0\), the positive saddle satisfies
\begin{equation}\label{eq:touchard-saddle-x}
  xr\ee^r=n,
  \qquad r=W_0(n/x).
\end{equation}
The phase derivatives are \(x\ee^r\Tch_j(r)\).  Since
\(x\ee^r=n/r\), the normalized quantities \(q_j(r)\) are exactly the same as
before.

\begin{theorem}[Universal Touchard expansion]\label{thm:touchard-universal}
Let \(K\subset(0,\infty)\) be compact.  Uniformly for \(x\in K\), with
\(r=W_0(n/x)\),
\begin{equation}\label{eq:touchard-allorders}
 \Bell_n(x)=
 \frac{n!\exp\!\bigl(x(\ee^r-1)\bigr)}
      {r^n\sqrt{2\pi n(1+r)}}
 \left(
   \sum_{m=0}^{M-1}\frac{c_m(r)}{n^m}
   +\BigO_{M,K}\!\left((r/n)^M\right)
 \right).
\end{equation}
The same \(c_m\) from \cref{eq:cm-partition} works for every \(x\).
After Stirling's expansion,
\begin{equation}\label{eq:touchard-log}
 \log\Bell_n(x)
 =x(\ee^r-1)+nr-n-\frac12\log(1+r)
 +\sum_{m=1}^{M-1}\frac{d_m(r)}{n^m}
 +\BigO_{M,K}((r/n)^M).
\end{equation}
\end{theorem}

\begin{proof}
Repeat \cref{eq:circular-integral} with
\(\ee^{r\ee^{\ii\theta}}-\ee^r\) multiplied by \(x\).  Equation
\cref{eq:touchard-saddle-x} gives the same quadratic coefficient
\(A=n(1+r)\) and the same scaled phase \cref{eq:scaled-phase}.  The
localization estimates are uniform for \(x\) in a compact positive interval
because \(r=W_0(n/x)\sim\log n\) uniformly there.  The logarithmic form follows
as in \cref{thm:log-W-form}.
\end{proof}

\begin{remark}
The universality is stronger than a mere substitution.  The parameter \(x\)
changes the saddle location and the leading exponential, but after the saddle
equation is imposed it disappears from every standardized local cumulant.
Regimes in which \(x=x_n\to0\) or \(x_n\to\infty\) require separate uniformity
conditions; \cref{thm:touchard-universal} intentionally treats fixed positive
parameter ranges.
\end{remark}

\subsection{A fixed-shift coefficient formula}

For fixed integer \(h\), one may either apply the main theorem directly to
\(N=n+h\), or keep the saddle \(r=W_0(n)\).  The latter gives
\begin{equation}\label{eq:shift-integral}
 \frac{\Bell_{n+h}}{(n+h)!}
 =\frac{\exp(\ee^r-1)}{2\pi r^{n+h}}
  \int_{-\pi}^{\pi}\exp(\Phi_r(\theta))\ee^{-\ii h\theta}\dd\theta.
\end{equation}
After the Gaussian scaling, define
\begin{equation}\label{eq:shift-cm}
 c_m(r;h):=
 \E\!\left[
  \coef{t}{2m}
  \exp\!\left(
   -\frac{\ii hZ}{\sqrt{1+r}}t
   +\sum_{s\ge1}a_s(Z;r)t^s
  \right)
 \right].
\end{equation}
Then, uniformly for \(|h|\le H\) with fixed \(H\),
\begin{equation}\label{eq:shift-expansion}
 \Bell_{n+h}=
 \frac{(n+h)!\exp(\ee^r-1)}{r^{n+h}\sqrt{2\pi n(1+r)}}
 \left(
   \sum_{m=0}^{M-1}\frac{c_m(r;h)}{n^m}
   +\BigO_{M,H}((r/n)^M)
 \right).
\end{equation}
Thus the fixed-shift coefficients are again finite complete Bell polynomials;
the only new atom is the linear Gaussian term in \cref{eq:shift-cm}.

\section{Algorithms and exact coefficient generation}\label{sec:algorithm}

The formulae above are not merely existential.  They give several independent
coefficient engines, each suited to a different task.  This section records the
engines in a form that can be implemented in an exact computer-algebra system
or formalized over the rational-function field \(\Q(r)\).

\subsection{A radical-free partition algorithm}

Although the standardized atoms \(q_j(r)\) in \cref{eq:qj-def} contain
half-integer powers of \(1+r\), every contributing monomial in
\cref{eq:cm-partition} has an even total Gaussian degree.  Define
\begin{equation}\label{eq:Uj-def}
 U_j(r):=\frac{\Tch_j(r)}{r}
        =\sum_{k=1}^{j}\Stwo{j}{k}r^{k-1}\in\N[r].
\end{equation}
For a vector \(\bm k=(k_1,\ldots,k_{2m})\) satisfying
\(\sum_{s=1}^{2m}s k_s=2m\), put \(K=\sum_s k_s\).  Since
\[
 \frac12\sum_{s=1}^{2m}(s+2)k_s=m+K,
\]
\cref{eq:cm-partition} becomes the completely radical-free identity
\begin{equation}\label{eq:cm-radical-free}
 \boxed{
 c_m(r)=
 \sum_{\substack{k_1,\ldots,k_{2m}\ge0\\
                   \sum_{s=1}^{2m}s k_s=2m}}
 \frac{(-1)^{m+K}(2m+2K-1)!!}{(1+r)^{m+K}}
 \prod_{s=1}^{2m}
 \frac{U_{s+2}(r)^{k_s}}{((s+2)!)^{k_s}k_s!}. }
\end{equation}
This is often the cleanest exact formula.  It requires one term for each
integer partition of \(2m\), not one term for each composition.  In
particular, it is finite before any simplification, and all of its arithmetic
takes place in \(\Q(r)\).

\begin{algorithm}[Partition engine for \(c_m\)]\label{alg:partition-engine}
Given \(m\ge0\):
\begin{enumerate}[label=\arabic*.,leftmargin=2.1em,itemsep=2pt]
 \item Generate \(\Tch_0=1\) and
       \(\Tch_{j+1}=r(\Tch_j+\Tch_j')\) up to \(j=2m+2\).
 \item Set \(U_j=\Tch_j/r\) for \(2\le j\le2m+2\).
 \item Enumerate the multiplicity vectors \(\bm k\) with
       \(\sum s k_s=2m\).
 \item Add the summand in \cref{eq:cm-radical-free} for each vector.
 \item Put the result over the common denominator \((1+r)^{3m}\), cancel
       common factors, and return the numerator polynomial.
\end{enumerate}
For \(m=0\), return \(c_0=1\) without entering the loop.
\end{algorithm}

The exponent \(m+K\) never exceeds \(3m\), because \(K\le2m\).  This makes the
denominator theorem \cref{thm:cm-structure} visible directly in the algorithm.
The degree estimate follows just as transparently: the product of the
\(U_{s+2}\)'s has degree at most
\(\sum(s+1)k_s=2m+K\), and multiplication by
\((1+r)^{3m-(m+K)}\) raises it to at most \(4m\).

\subsection{Dynamic Bell-polynomial algorithm}

Partition enumeration is conceptually direct, while the recurrence
\cref{eq:Pj-recurrence} is usually faster for moderate orders.  Let
\(\Gauss:\Q(r)[u]\to\Q(r)\) be the Gaussian moment functional
\begin{equation}\label{eq:gaussian-functional}
 \Gauss(u^{2j})=(2j-1)!!,
 \qquad
 \Gauss(u^{2j+1})=0.
\end{equation}
Construct \(a_s(u;r)\) from \cref{eq:scaled-phase}, then compute
\begin{equation}\label{eq:algorithm-P-recurrence}
 P_0(u;r)=1,
 \qquad
 P_j(u;r)=\frac1j\sum_{s=1}^{j}s\,a_s(u;r)P_{j-s}(u;r).
\end{equation}
The desired answer is \(c_m=\Gauss(P_{2m})\).  Only \(a_1,\ldots,a_{2m}\)
are needed.  Odd values of \(P_j\) need not be Gaussian-averaged, but they
must be retained because they feed later even values.

\begin{algbox}[title={Exact symbolic pseudocode}]
The following pseudocode uses the recurrence, but applies the Gaussian
functional before the final rational simplification.  The symbol
\texttt{stirling2(j,k)} denotes a Stirling number of the second kind.
\begin{lstlisting}[language=Python,basicstyle=\ttfamily\small,columns=fullflexible,
breaklines=true,showstringspaces=false]
def touchard(j, r):
    return sum(stirling2(j, k) * r**k for k in range(1, j + 1))

def gaussian_moment_polynomial(poly, u):
    ans = 0
    for exponent, coefficient in coefficients_in(poly, u):
        if exponent % 2 == 0:
            ans += coefficient * double_factorial(exponent - 1)
    return ans

def saddle_coefficient(m, r, u):
    if m == 0:
        return 1
    a = {}
    for s in range(1, 2*m + 1):
        T = touchard(s + 2, r)
        q = T / (r * (1 + r)**((s + 2)/2))
        a[s] = I**(s + 2) * q * u**(s + 2) / factorial(s + 2)
    P = {0: 1}
    for j in range(1, 2*m + 1):
        P[j] = expand(sum(s*a[s]*P[j-s] for s in range(1, j + 1))/j)
    return factor(gaussian_moment_polynomial(P[2*m], u))
\end{lstlisting}
In exact software, the exponent \((s+2)/2\) should be represented as a rational
number, not as floating-point division.  Alternatively, use
\cref{eq:cm-radical-free} and avoid radicals altogether.
\end{algbox}

\subsection{Logarithmic, exponential, and polynomial--logarithmic engines}

Once \(c_1,\ldots,c_M\) are known, no further integration is needed.
The remaining coefficient systems are ordinary formal-series transforms.
For convenient reference, the pipeline is
\begin{equation}\label{eq:coefficient-pipeline}
 \{c_m\}
 \xrightarrow{\ \log\ }
 \{\ell_m\}
 \xrightarrow{\ +\sigma_m\ }
 \{d_m\}
 \xrightarrow{\ d_m/r^m\ }
 \{\cD_m\}
 \xrightarrow{\ \exp\ }
 \{\cA_m\}.
\end{equation}
The arrows are implemented respectively by
\cref{eq:ell-recurrence,eq:d-m-def,eq:D-m-def,eq:A-recurrence}.  Explicitly,
\begin{align}
 \ell_m
 &=c_m-\frac1m\sum_{j=1}^{m-1}j\ell_jc_{m-j},
 \label{eq:algorithm-log}\\
 d_m&=\ell_m+\sigma_m,
 \qquad
 \cD_m=\frac{d_m}{r^m},
 \label{eq:algorithm-dD}\\
 \cA_0&=1,
 \qquad
 \cA_m=\frac1m\sum_{j=1}^{m}j\cD_j\cA_{m-j}.
 \label{eq:algorithm-exp}
\end{align}
All operations are triangular: coefficients through order \(M\) require only
coefficients through order \(M\) at the preceding stage.

For the polynomial--logarithmic sector, first compute the Lambert polynomials
\(\cR_N(\lambda)\) by \cref{eq:R-N-general}, then build the truncated formal
series
\begin{equation}\label{eq:algorithm-Delta}
 \Delta_M(z)=1-\lambda z+
             \sum_{j=1}^{M-1}\cR_j(\lambda)z^{j+1}.
\end{equation}
Ordinary truncated multiplication, reciprocal, and logarithm operations give
\(\cH_N\) from \cref{eq:H-general}, \(\cE_N\) from
\cref{eq:E-general}, and every \(p_{m,j}\) from \cref{eq:p-mj}.  Thus the
full transseries is generated by three familiar coefficient calculi:
Stirling-cycle numbers for \(W\), complete Bell polynomials for composition
and exponentiation, and Gaussian moments for the saddle integral.

\subsection{Complexity and internal checks}

Let \(p(N)\) be the integer partition function.  The direct formula for
\(c_m\) has \(p(2m)\) summands.  The dynamic recurrence has only
\(O(m^2)\) formal convolutions, although polynomial expansion and rational
simplification dominate the actual bit complexity.  In either approach, exact
coefficient generation is practical well beyond the four orders printed in
this article.

A reliable implementation should enforce the following invariants at every
order:
\begin{enumerate}[label=(\roman*),leftmargin=2.2em,itemsep=2pt]
 \item \((1+r)^{3m}c_m(r)\in\Q[r]\);
 \item its degree is at most \(4m\);
 \item \(c_m(r)=\BigO(r^m)\) as \(r\to\infty\);
 \item \(\cD_m(r)\) has a finite limit given by \cref{eq:D-limit};
 \item the limiting \(\cA_m\)'s satisfy the inverse Stirling identity
       \cref{eq:gamma-stirling-inverse};
 \item numerical substitution into \cref{eq:all-orders-main} improves the
       relative error by approximately one factor of \(r/n=\ee^{-r}\) per
       retained order.
\end{enumerate}
These checks detect sign errors, missing factorials, and incorrect Gaussian
moments very efficiently.

\section{Numerical validation}\label{sec:numerics}

\subsection{Definition of the tested approximants}

For \(M\ge0\), define the principal-saddle approximant
\begin{equation}\label{eq:SM-def}
 S_M(n):=
 \frac{n!\exp(\ee^r-1)}{r^n\sqrt{2\pi n(1+r)}}
 \sum_{m=0}^{M}\frac{c_m(r)}{n^m},
 \qquad r=W_0(n).
\end{equation}
Thus \(S_0\) is the classical leading approximation and \(S_4\) uses
\(c_1,c_2,c_3,c_4\), with \(c_4\) printed in
\cref{app:c4}.  Exact Bell numbers were generated by
\begin{equation}\label{eq:bell-recurrence-numerics}
 \Bell_0=1,
 \qquad
 \Bell_{n+1}=\sum_{k=0}^{n}\binom{n}{k}\Bell_k,
\end{equation}
and all transcendental evaluations were performed at substantially higher
precision than the digits shown.

\subsection{Relative-error table}

The table reports
\(\varepsilon_M(n):=S_M(n)/\Bell_n-1\).  The signs are retained because they
provide a more sensitive check than absolute errors alone.

\begin{table}[htbp]
\centering
\caption{Relative errors of the all-orders principal-saddle approximants.}
\label{tab:relative-errors}
\small
\setlength{\tabcolsep}{4.4pt}
\begin{tabular}{r@{\qquad}rrrrr}
\toprule
\(n\) & \(M=0\) & \(M=1\) & \(M=2\) & \(M=3\) & \(M=4\)\\
\midrule
10  & \( 2.6771612\,10^{-2}\) & \( 3.8253603\,10^{-4}\) & \(-1.2965341\,10^{-5}\) & \(-6.2268398\,10^{-6}\) & \(-1.1599281\,10^{-6}\)\\
20  & \( 1.5326828\,10^{-2}\) & \( 1.1293364\,10^{-4}\) & \(-3.8237417\,10^{-6}\) & \(-7.5548761\,10^{-7}\) & \(-7.1302490\,10^{-8}\)\\
50  & \( 7.2381178\,10^{-3}\) & \( 2.1197165\,10^{-5}\) & \(-5.8160607\,10^{-7}\) & \(-4.1043863\,10^{-8}\) & \(-1.3992028\,10^{-9}\)\\
100 & \( 4.0615707\,10^{-3}\) & \( 5.7437480\,10^{-6}\) & \(-1.2399466\,10^{-7}\) & \(-4.2053807\,10^{-9}\) & \(-6.2657058\,10^{-11}\)\\
200 & \( 2.2599385\,10^{-3}\) & \( 1.5057677\,10^{-6}\) & \(-2.4589871\,10^{-8}\) & \(-4.0769478\,10^{-10}\) & \(-2.4569360\,10^{-12}\)\\
500 & \( 1.0289535\,10^{-3}\) & \( 2.4364162\,10^{-7}\) & \(-2.6649303\,10^{-9}\) & \(-1.7364703\,10^{-11}\) & \(-2.4871987\,10^{-14}\)\\
\bottomrule
\end{tabular}
\end{table}

The data exhibit the predicted Poincar\'e ordering.  Because
\(r/n=\ee^{-r}\), the nominal ratio between consecutive terms is of order
\(r/n\), multiplied by a coefficient that grows polynomially in \(r\).  The
improvement is already substantial at \(n=10\), and by \(n=500\) four
corrections reduce the relative error to roughly \(2.5\cdot10^{-14}\).
The small nonmonotonicity between orders two and three at \(n=10\) is normal
for a finite asymptotic expansion and does not signal a failure of the
coefficient formula.

\subsection{Independent check of the first correction}

Using \(n=r\ee^r\), the logarithmic first-order coefficient gives
\begin{equation}\label{eq:first-check-expanded}
 \exp\!\left(\frac{d_1(r)}{n}+\BigO((r/n)^2)\right)
 =1-
 \frac{r^2(2r^2+7r+10)}{24(1+r)^3n}
 +\BigO\!\left(\frac{r^2}{n^2}\right).
\end{equation}
Consequently,
\begin{equation}\label{eq:first-check-final}
 \Bell_n=
 \frac{\exp\!\left(n(r-1+r^{-1})-1\right)}{\sqrt{1+r}}
 \left(
 1-
 \frac{r^2(2r^2+7r+10)}{24(1+r)^3n}
 +\BigO\!\left(\frac{r^2}{n^2}\right)
 \right),
\end{equation}
which is algebraically identical to the independently obtained 2026
first-correction formula discussed in \cref{sec:logtrans}.  This comparison is
important: it tests simultaneously the saddle coefficient \(c_1\), the
Stirling correction \(1/(12n)\), and the conversion from the factorial form to
the Lambert--\(W\) exponential form.

\section{Proof audit, transseries meaning, and limitations}\label{sec:audit}

\subsection{What has been proved}

The principal result is an ordinary asymptotic theorem with a moving
parameter.  For every fixed \(M\), the coefficients are evaluated at
\(r=W_0(n)\), and \cref{thm:all-orders} states a uniform remainder of relative
size \(\BigO_M((r/n)^M)\) after \(M\) retained coefficients.  Since
\(r\sim\log n\), \cref{cor:genuine-scale} proves that these terms form a
genuine asymptotic sequence.  The finite coefficient formulae are exact
identities in \(\Q(r)\), not asymptotic guesses.

The following implications are therefore rigorous:
\begin{equation}\label{eq:rigorous-implications}
 \text{localized Cauchy integral}
 \Longrightarrow \{c_m\}
 \Longrightarrow \{d_m,\cD_m,\cA_m\}
 \Longrightarrow \text{every finite truncation in }r.
\end{equation}
The polynomial--logarithmic formula then follows by composing any prescribed
finite truncation with the convergent large-positive-argument expansion of
\(W_0\), or formally to all orders by coefficient extraction.  At a fixed
requested precision, only finitely many \(W\)-coefficients and finitely many
saddle coefficients contribute.

\subsection{The ranking of the full transseries}

The phrase ``full transseries'' can be ambiguous unless the ordering is stated.
The expansion constructed here has three nested ranks:
\begin{enumerate}[label=\arabic*.,leftmargin=2.1em,itemsep=3pt]
 \item The \emph{extensive rank}, of size \(n\), is the expansion of
       \(n(r-1+r^{-1})\) in powers of \(L^{-1}\) with polynomial
       coefficients in \(\lambda=\log L\).
 \item The \emph{edge rank}, of size one, comes from
       \(-1-\tfrac12\log(1+r)\) and has its own
       \(L^{-1}\Q[\lambda]\)-series.
 \item The \emph{exponentially small ranks} are indexed by \(m\ge1\) and have
       magnitude
       \(\ee^{-mr}=n^{-m}L^m(1+o(1))\), each carrying a complete
       \(L^{-1}\Q[\lambda]\)-expansion.
\end{enumerate}
Within a fixed exponential sector, powers of \(L^{-1}\) give the local
ordering.  Across sectors, every term of the \(m+1\) sector is beyond all
algebraic orders of the \(m\) sector after a finite polynomial--logarithmic
truncation.  This is the precise sense in which
\cref{eq:full-polylog-box} is a transseries rather than one undifferentiated
formal series.

\subsection{What is not claimed}

\begin{auditbox}[title={Boundaries of the present result}]
\begin{enumerate}[label=(\alph*),leftmargin=2.1em,itemsep=3pt]
 \item No convergence claim is made for the infinite principal-saddle series
       \(\sum c_m(r)n^{-m}\) when the order tends to infinity together with
       \(n\).  The theorem is Poincar\'e asymptotic: \(M\) is fixed first.
 \item No sharp factorial-over-power estimate for \(c_m\) as
       \(m\to\infty\) is proved.  Such an estimate is needed for rigorous
       optimal truncation and resurgence.
 \item The formal nonprincipal saddle germs in \cref{sec:complex} do not by
       themselves determine which complex saddles occur on a deformation of
       the original Cauchy circle, nor their Stokes multipliers.
 \item The Touchard theorem is uniform only for \(x\) in a fixed compact
       subset of \((0,\infty)\).  Coalescing regimes such as \(x\to0\), or
       complex \(x\) near Stokes curves, require a separate analysis.
 \item The document supplies a formalization map but not a machine-checked
       proof.  In particular, the complex-analytic localization estimates and
       all truncation bookkeeping would have to be encoded explicitly in a
       proof assistant.
\end{enumerate}
\end{auditbox}

These qualifications are mathematically substantive.  They separate the
complete all-orders coefficient calculus proved here from a stronger, still
open hyperasymptotic description of the exact Bell sequence.

\subsection{A formalization map}

The proof naturally decomposes into modules.
\begin{longtable}{p{0.25\textwidth}p{0.66\textwidth}}
\toprule
Module & Principal statements and dependencies\\
\midrule
\endfirsthead
\toprule
Module & Principal statements and dependencies\\
\midrule
\endhead
Combinatorial algebra & Stirling recurrence; Touchard recurrence
\cref{eq:touchard-recurrence}; divisibility \(r\mid\Tch_j\); complete Bell
polynomials; integer-partition indexing.\\[2pt]
Formal power series & Exponential recurrence \cref{eq:Pj-recurrence}; logarithm
recurrence \cref{eq:ell-recurrence}; exponential recurrence
\cref{eq:A-recurrence}; truncated composition and inversion.\\[2pt]
Gaussian calculus & Moment identity \cref{eq:gaussian-functional}; Wick pairing;
equivalence of \cref{eq:cm-bellpoly,eq:cm-partition}.\\[2pt]
Real asymptotics & Bounds for \(r=W_0(n)\); estimates
\cref{eq:central-decay,eq:outer-decay}; domination of Taylor remainders by a
Gaussian integrable majorant.\\[2pt]
Complex analysis & Cauchy's coefficient formula; parameterization of the
saddle circle; complementary-arc localization.\\[2pt]
Lambert calculus & Stirling-cycle-number expansion
\cref{eq:R-N-general}; finite coefficient extraction for
\(\cH_N,\cE_N,p_{m,j}\).\\
\bottomrule
\end{longtable}
The algebraic modules can be developed independently of analysis and reused
for other admissible exponential generating functions.  The only Bell-specific
analytic input is the phase \(\ee^{r\ee^{\ii\theta}}\), whose derivatives are
encoded by the Touchard polynomials.

\section{Further deductions and research directions}\label{sec:future}

The coefficient calculus reveals several questions that are not visible at
leading order.

\subsection{Large-order growth and resurgence}

The nearest competing saddles are \(W_{\pm1}(n)\).  Their action gap is of
magnitude \(2\pi^2n/r^2\) by \cref{eq:real-action-gap}.  Standard resurgence
heuristics therefore suggest that the late terms of the principal-saddle
series should eventually be controlled by these adjacent actions.  A natural
program is to prove a bound of schematic form
\begin{equation}\label{eq:late-term-program}
 c_m(r)\sim
 \sum_{k=\pm1}
 \frac{S_k(r)\,\Gamma(m+\alpha_k)}{\Delta_k(r)^{m+\alpha_k}}
 \left(1+\frac{b_{k,1}(r)}m+\cdots\right),
\end{equation}
with explicitly determined Stokes constants \(S_k\).  Formula
\cref{eq:cm-partition} supplies arbitrarily many exact low-order data with
which such a conjecture can be tested, but it does not prove
\cref{eq:late-term-program}.

\subsection{Ratio, logarithmic derivative, and finite differences}

Because all transforms are formal and triangular, the present results
immediately generate complete expansions of
\begin{equation}\label{eq:derived-observables}
 \frac{\Bell_{n+1}}{\Bell_n},
 \qquad
 \log\Bell_{n+1}-\log\Bell_n,
 \qquad
 \Delta^j\log\Bell_n,
 \qquad
 \frac{\Bell_{n+h}}{\Bell_n}
\end{equation}
for fixed \(j,h\).  There are two equivalent routes: substitute
\(n\mapsto n+h\) into the master transseries and re-expand, or use the shifted
coefficient germ \cref{eq:shift-cm}.  Comparing the two gives an infinite
family of exact identities among the rational functions \(c_m(r;h)\), the
derivatives of \(W\), and the unshifted \(c_m(r)\).  These identities appear
well suited to automated proof and can sharpen explicit ratio bounds.

\subsection{Uniform Touchard regimes}

The universality theorem suggests several boundary problems:
\begin{enumerate}[label=(\roman*),leftmargin=2.2em,itemsep=2pt]
 \item \(x=x_n\to0^+\), where \(r=W(n/x)\) grows faster than \(\log n\);
 \item \(x=x_n\to\infty\), where the saddle can leave the logarithmic regime;
 \item negative or complex \(x\), where several saddles can have comparable
       real action;
 \item double scaling near a Stokes curve, requiring an error-function or
       Airy-type smoothing rather than a single Gaussian saddle.
\end{enumerate}
The finite local coefficients remain meaningful by analytic continuation, but
the global contour and uniform remainder must be rebuilt in each regime.

\subsection{A general exponential-of-exponential template}

Consider a family
\begin{equation}\label{eq:general-template}
 F(z)=\exp(H(z)),
 \qquad
 [z^n]F(z)=\frac{F(r)}{r^n\sqrt{2\pi b(r)}}
           \left(1+\cdots\right),
\end{equation}
where \(rH'(r)=n\) and
\(b(r)=(r\partial_r)^2H(r)\).  Define normalized logarithmic derivatives
\begin{equation}\label{eq:general-cumulants}
 \kappa_j(r):=\frac{(r\partial_r)^jH(r)}{b(r)^{j/2}}.
\end{equation}
Exactly the same Gaussian Bell-polynomial calculation gives every local
saddle coefficient as a finite polynomial in
\(\kappa_3,\ldots,\kappa_{2m+2}\).  Bell numbers are special because
\(H(z)=\ee^z-1\) converts these derivatives into Touchard polynomials and the
saddle into a Lambert \(W\)-function.  This observation provides a systematic
route from the repository's combinatorial coefficient calculus to a wider
class of labelled combinatorial assemblies.

\section{Conclusion}\label{sec:conclusion}

The Bell-number saddle is classical; the point of this article is to make its
entire correction mechanism explicit and finite.  The result can be summarized
in four layers.

First, all local derivatives of the phase are Touchard polynomials.  Second,
Gaussian reduction turns those derivatives into the exact finite coefficient
formula \cref{eq:cm-partition}, or equivalently the radical-free form
\cref{eq:cm-radical-free}.  Third, ordinary logarithm and exponential
recurrences combine the saddle series with Stirling's series to produce the
bounded rational coefficient families \(\cD_m(r)\) and \(\cA_m(r)\) in the
natural exponential variable \(\ee^{-r}=r/n\).  Fourth, the
Stirling-cycle-number expansion of \(W_0(n)\) resolves every one of these
coefficients into the ranked polynomial--logarithmic transseries in
\(L=\log n\) and \(\lambda=\log L\).

Thus every coefficient at every displayed rank is governed by a finite
formula.  The theorem includes an arbitrary-order principal-saddle remainder,
its first correction agrees with an independent recent derivation, and direct
numerics confirm rapid improvement.  The remaining frontier is no longer the
local coefficient calculus; it is the global hyperasymptotic problem of
nonprincipal saddles, large-order coefficient growth, and Stokes switching.

\appendix

\section{The fourth principal-saddle coefficient}\label{app:c4}

For numerical checks and as a higher-order checksum, the fourth coefficient is
\begin{equation}\label{eq:c4-explicit}
 c_4(r)=-\frac{N_4(r)}{39813120(1+r)^{12}},
\end{equation}
where
\begin{align*}
 N_4(r)={}&
 9136r^{16}-357408r^{15}-5092600r^{14}-29432904r^{13}
 -98499653r^{12}\\
 &-212369184r^{11}-302729704r^{10}-277448328r^9
 -156972248r^8\\
 &-52300272r^7+110053352r^6+302787360r^5+198352r^4\\
 &+2545056r^3+647456r^2+109632r+9136.
\end{align*}
The sign in \cref{eq:c4-explicit} applies to the whole numerator.  In
particular, the leading behavior is
\begin{equation}\label{eq:c4-leading}
 c_4(r)=-\frac{571}{2488320}r^4+\BigO(r^3),
\end{equation}
in agreement with \(\gamma_4=-571/2488320\) in
\cref{eq:gamma-values}.

\section{Coefficient dictionaries}\label{app:dictionaries}

This appendix collects the coefficient-extraction rules in one place.  They
are useful when translating between notational conventions in the literature.

\subsection{Principal saddle}

With \(Z\sim N(0,1)\),
\begin{align}
 a_s(Z;r)
 &=\frac{\ii^{s+2}\Tch_{s+2}(r)}
         {(s+2)!\,r(1+r)^{(s+2)/2}}Z^{s+2},
 \label{eq:dict-a}\\
 c_m(r)
 &=\E\coef{t}{2m}
   \exp\!\left(\sum_{s\ge1}a_s(Z;r)t^s\right).
 \label{eq:dict-c}
\end{align}
Equations \cref{eq:dict-a,eq:dict-c} are the shortest universal definition;
\cref{eq:cm-radical-free} is the shortest radical-free definition.

\subsection{Logarithmic and multiplicative sectors}

Let \(C(x)=1+\sum_{m\ge1}c_mx^m\).  Then
\begin{align}
 \log C(x)&=\sum_{m\ge1}\ell_mx^m,
 &d_m&=\ell_m+\sigma_m,
 &\cD_m&=r^{-m}d_m,
 \label{eq:dict-log}\\
 \exp\!\left(\sum_{m\ge1}\cD_my^m\right)
 &=\sum_{m\ge0}\cA_my^m.
 \label{eq:dict-mult}
\end{align}
The Bell number can therefore be written either as
\begin{align}
 \log\Bell_n
 &\sim n(r-1+r^{-1})-1-\tfrac12\log(1+r)
       +\sum_{m\ge1}\cD_m(r)\ee^{-mr},
 \label{eq:dict-logbell}\\
 \Bell_n
 &\sim \frac{\exp(n(r-1+r^{-1})-1)}{\sqrt{1+r}}
       \sum_{m\ge0}\cA_m(r)\ee^{-mr}.
 \label{eq:dict-multbell}
\end{align}

\subsection{Inverse-\texorpdfstring{\(r\)}{r} and polynomial--logarithmic subcoefficients}

Because every \(\cD_m(r)\) is regular at \(r=\infty\), define
\begin{equation}\label{eq:Dmk-def}
 \cD_{m,j}:=\coef{u}{j}\cD_m(u^{-1}),
 \qquad
 \cD_m(r)=\sum_{j\ge0}\cD_{m,j}r^{-j}
\end{equation}
formally at infinity.  Likewise
\(\cA_{m,j}:=\coef{u}{j}\cA_m(u^{-1})\).  Substitution of
\(r=L\Delta(L^{-1})\) gives the sector coefficients
\begin{equation}\label{eq:sector-double-extraction}
 p_{m,j}(\lambda)
 =\coef{z}{j}
   \left[z^m d_m\!\left(z^{-1}\Delta(z)\right)\right]
 =\coef{z}{j}
   \left[\Delta(z)^m\cD_m\!\left(z^{-1}\Delta(z)\right)\right].
\end{equation}
This second equality makes explicit that the leading value is controlled by
\(\cD_m(\infty)\), whereas all lower inverse powers of \(r\) generate the
internal \(L^{-1}\)-tail of the same exponential sector.

\section{A direct derivation of the first coefficient}\label{app:c1proof}

For readers wishing to audit every normalization, retain terms through order
\(n^{-1}\) in \cref{eq:scaled-phase}:
\begin{equation}\label{eq:first-local-phase}
 \exp\!\left(n^{-1/2}a_1+n^{-1}a_2+\BigO(n^{-3/2})\right)
 =1+n^{-1/2}a_1+n^{-1}\left(a_2+\frac12a_1^2\right)
  +\BigO(n^{-3/2}).
\end{equation}
The odd term has zero Gaussian mean.  Since
\[
 a_1=-\frac{\ii q_3}{6}Z^3,
 \qquad
 a_2=\frac{q_4}{24}Z^4,
\]
we obtain
\begin{align}
 c_1
 &=\frac{q_4}{24}\E Z^4
   -\frac{q_3^2}{72}\E Z^6
 =\frac{q_4}{8}-\frac{5q_3^2}{24}.
 \label{eq:c1-audit-q}
\end{align}
Now
\begin{align*}
 \Tch_3(r)&=r+3r^2+r^3,
 &\Tch_4(r)&=r+7r^2+6r^3+r^4,
\end{align*}
so substitution and collection over \(24(1+r)^3\) yield
\[
 c_1(r)=-\frac{2r^4+9r^3+16r^2+6r+2}{24(1+r)^3}.
\]
Finally, adding the first Stirling logarithm coefficient gives
\begin{align*}
 d_1(r)
 &=c_1(r)+\frac1{12}\\
 &=-\frac{r^2(2r^2+7r+10)}{24(1+r)^3},
\end{align*}
which explains the cancellation of the constant and linear terms in the
numerator.  This cancellation is the first visible instance of the Bernoulli
boundary mechanism in \cref{prop:D-limits}.

\section{Bibliographic notes}\label{app:bibnotes}

Moser and Wyman identified the Lambert-type saddle and developed classical
correction expansions.  Hayman's admissibility theory places the Bell
generating function in a broad coefficient-asymptotic framework, while de
Bruijn gives the familiar logarithmic form.  Harris and Schoenfeld developed
general coefficient expansions for analytic functions.  Corless et al. and
the DLMF provide the branch-aware all-orders Lambert \(W\) expansion used in
\cref{sec:polylog}.  More recent work supplies explicit bounds and an
independent elementary derivation of the first Bell correction.

The novelty claimed here is organizational and algebraic rather than a claim
that saddle-point expansions for Bell numbers were previously unknown: the
article makes every coefficient finite and computable, proves the rational
structure adapted to \(r=W(n)\), and threads the result through logarithmic,
exponential, and polynomial--logarithmic transseries in one consistent
notation.

\begin{thebibliography}{99}

\bibitem{Bell1934}
E.~T. Bell,
\newblock Exponential polynomials,
\newblock \emph{Annals of Mathematics} \textbf{35} (1934), 258--277.

\bibitem{Comtet1974}
L.~Comtet,
\newblock \emph{Advanced Combinatorics: The Art of Finite and Infinite
Expansions},
\newblock D.~Reidel, Dordrecht, 1974.

\bibitem{MoserWyman1955}
L.~Moser and M.~Wyman,
\newblock An asymptotic formula for the Bell numbers,
\newblock \emph{Transactions of the Royal Society of Canada, Section III}
\textbf{49} (1955), 49--54.

\bibitem{Hayman1956}
W.~K. Hayman,
\newblock A generalisation of Stirling's formula,
\newblock \emph{Journal f\"ur die reine und angewandte Mathematik}
\textbf{196} (1956), 67--95.

\bibitem{HarrisSchoenfeld1968}
B.~Harris and L.~Schoenfeld,
\newblock Asymptotic expansions for the coefficients of analytic functions,
\newblock \emph{Illinois Journal of Mathematics} \textbf{12} (1968),
264--277.

\bibitem{deBruijn1981}
N.~G. de~Bruijn,
\newblock \emph{Asymptotic Methods in Analysis}, third edition,
\newblock Dover, New York, 1981.

\bibitem{CorlessEtAl1996}
R.~M. Corless, G.~H. Gonnet, D.~E.~G. Hare, D.~J. Jeffrey, and D.~E. Knuth,
\newblock On the Lambert \(W\) function,
\newblock \emph{Advances in Computational Mathematics} \textbf{5} (1996),
329--359.

\bibitem{FlajoletSedgewick2009}
P.~Flajolet and R.~Sedgewick,
\newblock \emph{Analytic Combinatorics},
\newblock Cambridge University Press, Cambridge, 2009.

\bibitem{Hwang2023}
H.-K. Hwang,
\newblock A curious identity arising from Stirling's formula and saddle-point
method on two different contours,
\newblock \emph{The Electronic Journal of Combinatorics} \textbf{30}(4)
(2023), Paper P4.3.

\bibitem{GrunwaldSerafin2025}
J.~Grunwald and G.~Serafin,
\newblock Explicit bounds for Bell numbers and their ratios,
\newblock \emph{Journal of Mathematical Analysis and Applications}
\textbf{549}(2) (2025), Article 129527,
\newblock doi:10.1016/j.jmaa.2025.129527.

\bibitem{HwangLiZacharovas2026}
H.-K. Hwang, C.-Y. Li, and V.~Zacharovas,
\newblock Elementary asymptotics for the Stirling numbers of the second kind:
the central range,
\newblock arXiv:2605.29633, 2026.

\bibitem{DLMFW}
NIST Digital Library of Mathematical Functions,
\newblock \S4.13, Lambert \(W\)-function,
\newblock release consulted September 2026.

\bibitem{RepoCoefficientCalculus}
V.~Reshetnikov,
\newblock \emph{Combinatorial Coefficient Calculus},
\newblock \href{https://github.com/VladimirReshetnikov/ProveIt/tree/main/Analysis/FabiusFunction/docs/semi-formalized-research-frontiers/drafts/combinatorial-coefficient-calculus/Combinatorial_Coefficient_Calculus}{repository manuscript in the ProveIt project},
revision consulted September 3, 2026.

\bibitem{RepoSeriesTransseries}
V.~Reshetnikov,
\newblock \emph{Series and Transseries},
\newblock \href{https://github.com/VladimirReshetnikov/ProveIt/tree/main/Analysis/FabiusFunction/docs/semi-formalized-research-frontiers/drafts/series-and-transseries}{repository draft collection in the ProveIt project},
revision consulted September 3, 2026.

\end{thebibliography}

\end{document}
